\documentclass[11pt]{amsart}
\usepackage{amsmath,amssymb,amsthm,mathtools}
\usepackage[margin=1.15in]{geometry}
\usepackage{xcolor}
\usepackage[colorlinks=true,linkcolor=blue!60!black,citecolor=blue!60!black,urlcolor=blue!60!black]{hyperref}

\newtheorem{theorem}{Theorem}
\renewcommand{\thetheorem}{\Alph{theorem}}
\newtheorem*{bksthm}{Theorem 1.3 of \cite{BKS}}
\newtheorem{lemma}{Lemma}
\newtheorem{proposition}[lemma]{Proposition}
\newtheorem{corollary}[lemma]{Corollary}
\theoremstyle{remark}
\newtheorem{remark}[lemma]{Remark}

\newcommand{\R}{\mathbb{R}}
\newcommand{\Z}{\mathbb{Z}}
\newcommand{\dd}{\,\mathrm{d}}
\newcommand{\abs}[1]{\lvert #1\rvert}
\newcommand{\norm}[1]{\lvert #1\rvert}
\newcommand{\supn}[1]{\lvert #1\rvert_\infty}
\newcommand{\1}{\mathbf{1}}
\newcommand{\Qd}{\mathcal{Q}_\delta}

\title[A gap in Bux--Kassmann--Schulze and a repair]{The discretization step in
Bux--Kassmann--Schulze:\\ a gap and a repair by random lattice sampling}
\author{Claude}
\address{Anthropic}
\date{September 2026}



\begin{document}
\begin{abstract}
Bux, Kassmann and Schulze prove that quadratic forms whose kernels are nondegenerate only on
double cones of arbitrary orientation are comparable, on every ball, to the Gagliardo seminorm of
$H^{\alpha/2}$. Their proof passes from a discrete comparability theorem on $\Z^d$ to the
continuous one by a limit which, as written, requires the function to lie a priori in
$H^{\alpha/2}$ of a larger ball, which is what is to be proved. We explain the gap and close
it along the paper's own route. Instead of averaging the function over lattice cells, we sample
it at the points of a randomly shifted and randomly distorted lattice, apply the discrete theorem
to each lattice, and average over the lattices. This proves the enlarged-ball comparability for
every measurable function, with no finiteness assumption, and with it all the continuous results
of the paper. The argument has been checked in Lean~4.
\end{abstract}
\maketitle

\section{Introduction}\label{sec:intro}

We use the notation of Bux, Kassmann and Schulze~\cite{BKS}, and cite the numbering of the arXiv
version~\cite{BKS}. A \emph{double cone} with axis a unit vector $v\in\R^d$ and apex angle
$\theta\in(0,\pi/2]$ is
\[
  V=\{\,u\in\R^d\setminus\{0\} : \abs{\langle u,v\rangle}>\cos\theta\,\norm{u}\,\},
\]
the set of nonzero vectors making an angle less than $\theta$ with the line $\R v$
\cite[Definition~2.1]{BKS}. A \emph{configuration} $\Gamma$ assigns to each $x\in\R^d$ a double
cone $\Gamma(x)$, and $V^\Gamma[x]=x+\Gamma(x)$. It is \emph{$\vartheta$-bounded} if the infimum
$\vartheta$ of the apex angles of the cones in $\Gamma(\R^d)$ is positive. Let $\alpha\in(0,2)$ and $\Lambda\ge1$. We
consider measurable $k\colon\R^d\times\R^d\to[0,\infty]$ with $k(x,y)=k(y,x)$ and
\begin{equation}\label{eq:2}
  \Lambda^{-1}\bigl(\1_{V^\Gamma[x]}(y)+\1_{V^\Gamma[y]}(x)\bigr)\norm{x-y}^{-d-\alpha}
  \le k(x,y)\le \Lambda\norm{x-y}^{-d-\alpha}\qquad(x,y\in\R^d).\tag{2}
\end{equation}
This is condition (2) of \cite{BKS}. For a set $D\subseteq\R^d$ and $f\colon\R^d\to\R$ write
\[
  \abs{f}^2_{H_k(D)}=\iint_{D\times D}\bigl(f(y)-f(x)\bigr)^2k(x,y)\dd x\dd y,\qquad
  \abs{f}^2_{H^{\alpha/2}(D)}=\iint_{D\times D}\frac{(f(y)-f(x))^2}{\norm{x-y}^{d+\alpha}}\dd x\dd y,
\]
with values in $[0,\infty]$.

The main result of \cite{BKS}, Theorem~1.1, is that for a $\vartheta$-admissible $\Gamma$ (see
Section~\ref{sec:consequences}) there is $c$, depending only on $d$, $\vartheta$, $\Lambda$ and
$\alpha$, such that $\abs{f}^2_{H^{\alpha/2}(B)}\le c\,\abs{f}^2_{H_k(B)}$ for every ball $B$ and
every $f\in L^2(B)$. Moreover, for each $\alpha_0\in(0,2)$, $c$ can be chosen the same for all
$\alpha\in[\alpha_0,2)$. The proof has three steps.
\begin{enumerate}
\item A discrete analogue on $\Z^d$, Theorem~1.3 of \cite{BKS}, restated in Section~3.1. It is the heart of
  the paper.
\item The passage from the discrete to the continuous setting, in \S3.2 of \cite{BKS} (proof of
  Lemma~3.7). It yields the \emph{enlarged-ball} inequality
  $c\,\abs{f}^2_{H^{\alpha/2}(B)}\le\abs{f}^2_{H_k(B^*)}$, with $B=B_R(x_0)$ and $B^*=B_{\kappa R}(x_0)$.
\item The removal of the enlargement, Lemma~A.1 of \cite{BKS}, via a Whitney decomposition and an
  inequality of Dyda~\cite{Dyda}.
\end{enumerate}

Step~2 contains a gap, described in Section~\ref{sec:gap}. As written, the limit $h\to0$ on the
kernel side of the discretized inequality is justified by dominated convergence, but no dominating
function is available for a general $f\in H_k(B^*)$, and the lemma is applied, and must be
applied, to every such $f$. Separately, for $\alpha>1$ the limit asserted there is false: the
discretization mishandles the cells at the boundary of $B^*$ (Section~\ref{sec:gap}).

Our main result closes the gap. It uses nothing from \cite{BKS} beyond Theorem~1.3.

\begin{theorem}\label{thm:A}
Let $d\ge1$, $\vartheta\in(0,\pi/2]$, $\Lambda\ge1$ and $\alpha\in(0,2)$. There are $\kappa\ge1$
and $C<\infty$, depending only on $d$, $\vartheta$, $\Lambda$ and $\alpha$, with the following
property. Let $\Gamma$ be a $\vartheta$-bounded configuration, which need not satisfy any
measurability condition. Let $k$ be measurable, symmetric and satisfy \eqref{eq:2}, and let
$f\colon\R^d\to\R$ be measurable. Then for every ball $B=B_R(x_0)$,
\[
  \abs{f}^2_{H^{\alpha/2}(B_R(x_0))}\le C\,\abs{f}^2_{H_k(B_{\kappa R}(x_0))}.
\]
\end{theorem}

No finiteness of either side is assumed. This is exactly the enlarged-ball inequality that
Step~2 is meant to provide, now for every measurable $f$, so Steps~1 and~3 complete the proof of
Theorem~1.1 of \cite{BKS} (Section~\ref{sec:consequences}).

\medskip\noindent\textbf{The idea.} In \cite{BKS} the function is replaced by its averages
$f_h$ over the cells of $h\Z^d$. The averages destroy the oscillation of $f$ inside a cell, and
recovering the continuous form in the limit is what requires domination. We instead
\emph{sample} $f$ at the points of a lattice. Sampling keeps every difference $f(s)-f(t)$ between
lattice points exactly, so the discrete theorem applies to the sampled function for every
lattice, with no limit to take. The lattice is then made random, by a shift
$\eta\in[-\tfrac12,\tfrac12)^d$ and a small linear distortion $A=I+E$. Averaging the discrete
inequality over the lattices reproduces both continuous forms, up to constants and up to pairs at
distance $O(h)$. The random shift alone would sample only differences in directions of $h\Z^d$;
the random distortion spreads the differences over all directions.

\medskip\noindent\textbf{Related work.} Chaker and Silvestre \cite[Theorem~1.3]{CS} prove
$\iint_{B_2\times B_2}(u(x)-u(y))^2K\ge c\lambda\abs{u}^2_{\dot H^s(B_1)}$ for every $u$, under
their Assumption~1.1: for every ball $B$ and $x\in B$, the set of $z\in B$ with
$K(x,z)\ge\lambda\norm{x-z}^{-d-2s}$ has measure at least $\mu\abs B$. They remark that their
result implies that of \cite{BKS}. A kernel satisfying \eqref{eq:2} does satisfy Assumption~1.1,
with $s=\alpha/2$, $\lambda=\Lambda^{-1}$ and $\mu=(\sin^2\vartheta/16)^d$ (Lemma~\ref{lem:cs}), so
their theorem applies (Remark~\ref{rem:cs}). Their method is different: an iteration on kernels with a covering
argument, not a discretization. So, with Lemma~A.1 of \cite{BKS}, the continuous conclusions of
\cite{BKS} were already available through \cite{CS}. The point of the present note is that the
discretization route of \cite{BKS}, through their Theorem~1.3, can itself be completed. We are
not aware of any earlier discussion of the gap.

\medskip\noindent\textbf{Formal verification.} Theorem~\ref{thm:A} and every lemma below have
been formalized and checked in Lean~4 with Mathlib (Section~\ref{sec:lean}). The formalization
also proves Theorem~1.3 of \cite{BKS} itself.

\section{The gap}\label{sec:gap}

We follow the proof of Lemma~3.7 of \cite{BKS} (pp.~14--16). Let $f\in H_k(B^*)$. For $h\in(0,1)$
let $f_h(x)=h^{-d}\int_{A_h(x)\cap B^*}f$ for $x\in h\Z^d\cap B^*$, where $A_h(x)$ is the cell of
side $h$ about $x$. By their Corollary~3.6 there are a constant $C$ and a $\vartheta'$-bounded
configuration $\Gamma_h$, for some $\vartheta'>0$, for which the discretized kernel $\omega^k_h$
satisfies the discrete bounds (8) of \cite{BKS} on $h\Z^d$, with $R_0=\sqrt d$ and $\Lambda=C$. Their
Corollary~3.1 (Theorem~1.3 on $h\Z^d$) and Lemma~3.4 then give the discrete inequality (15) of
\cite{BKS}:
\begin{multline*}
  c\sum_{\substack{x,y\in B\cap h\Z^d\\ \abs{x-y}>\sqrt d h}}(f_h(x)-f_h(y))^2
    \int_{A_h(x)\times A_h(y)}\abs{s-t}^{-d-\alpha}\dd(s,t)\\
  \le\sum_{\substack{x,y\in B^*\cap h\Z^d\\ \abs{x-y}>\sqrt d h}}(f_h(x)-f_h(y))^2
    \int_{A_h(x)\times A_h(y)}k(s,t)\dd(s,t).
\end{multline*}
The right-hand side is $\int g_h$, where $g_h(s,t)=(f_h(x)-f_h(y))^2k(s,t)$ for
$(s,t)\in\tilde A_h(x)\times\tilde A_h(y)$ with $x,y$ as in the sum, and $g_h=0$ otherwise.
By Lebesgue differentiation (their Lemma~A.2), $g_h\to g=(f(s)-f(t))^2k(s,t)\1_{B^*\times B^*}$
almost everywhere. The proof then states, on p.~16:
\begin{quote}
``For the right hand side in (15) this implies with help of dominated convergence
[$\,\int g_h\to\int g\,$].''
\end{quote}
No dominating function is given, and none is available from the hypothesis $f\in H_k(B^*)$.

Moreover, as printed the step is false when $\alpha>1$, for a reason unrelated to domination. For
a cell that meets $\partial B^*$, $f_h(x)$ is still normalized by $h^{-d}$, so it is not an average
of $f$. Take $d=1$, $B^*=(-1,1)$, $f\equiv1$ and $k(x,y)=\abs{x-y}^{-1-\alpha}$, which satisfies
\eqref{eq:2} with $\Lambda=2$ for every $\Gamma$. Then $g=0$. For $h=1/(n+\frac14)$ the points of
$h\Z\cap B^*$ are $jh$, $\abs j\le n$, and the cell of $nh=1-h/4$ gives $f_h(nh)=\frac34$, while
$f_h=1$ at the interior points. For $s\in A_h(nh)$ and $t\in A_h(nh-2h)$ we have
$\abs{s-t}\le3h$, so this single pair of points contributes to $\int g_h$ at least
\[
  \frac1{16}\,h^2\,(3h)^{-1-\alpha}=\frac{h^{1-\alpha}}{16\cdot3^{1+\alpha}},
\]
which tends to infinity as $h\to0$ (formalized as \texttt{QFS.BoundaryDefect.rhs15\_tendsto\_top}).
This defect goes away if $f_h$ is formed on a slightly larger
ball, so that every cell used lies inside it. The domination problem does not.

For cells inside $B^*$, the natural upper bound on $g_h$ comes from Jensen's inequality,
\[
  (f_h(x)-f_h(y))^2\le h^{-2d}\int_{A_h(x)\times A_h(y)}(f(s')-f(t'))^2\dd(s',t'),
\]
together with the upper bound in \eqref{eq:2},
$k(s,t)\le\Lambda\norm{s-t}^{-d-\alpha}\lesssim\norm{s'-t'}^{-d-\alpha}$.
This bounds $g_h$ by local averages of $(f(s')-f(t'))^2\norm{s'-t'}^{-d-\alpha}$. Any integrable
majorant built this way requires at least $f\in H^{\alpha/2}(B^*)$, which is the conclusion
being proved.

The hypothesis controls $(f(s)-f(t))^2k(s,t)$, and this cannot be used here. The function $g_h$
pairs the value $k(s,t)$ with differences of \emph{averages} of $f$ over the two cells, and the
lower bound in \eqref{eq:2} controls $k$ only on the cones. The kernel may vanish at $(s',t')$
while being of full size at $(s,t)$ in the same pair of cells.

Nor can the step be rescued by approximating $f\in H_k(B^*)$ by smooth functions in the $H_k$
norm: that would require the density statement of Theorem~1.4, whose proof uses Lemma~3.7.
The final inequality of Lemma~3.7 is true, by Theorem~\ref{thm:A}, or by \cite{CS} through
Lemma~\ref{lem:cs}. But its proof
in \cite{BKS} is incomplete at this step.

\section{Proof of Theorem~\ref{thm:A}}\label{sec:proof}

Throughout, $d\ge1$, $\norm{\cdot}$ is the Euclidean norm and $\supn{\cdot}$ the maximum norm, and
$\abs{\,\cdot\,}$ also denotes Lebesgue measure. Balls $B_r(x)$ are open and $\bar B_r(x)$
closed. Integrals of $[0,\infty]$-valued functions are Lebesgue integrals, and all identities
and inequalities are between elements of $[0,\infty]$. We use the convention $0\cdot\infty=0$, and read
the integrand of $\abs{f}^2_{H^{\alpha/2}}$ as $0$ on the diagonal.

\subsection{The discrete theorem}\label{sec:discrete}

\begin{bksthm}
Let $\Gamma$ be a $\vartheta$-bounded configuration and $\alpha\in(0,2)$. Let
$\omega\colon\Z^d\times\Z^d\to[0,\infty]$ satisfy $\omega(x,y)=\omega(y,x)$ and
\[
  \Lambda^{-1}\bigl(\1_{V^\Gamma[x]}(y)+\1_{V^\Gamma[y]}(x)\bigr)\abs{x-y}^{-d-\alpha}
  \le\omega(x,y)\le\Lambda\abs{x-y}^{-d-\alpha}
\]
for $\abs{x-y}>R_0$, where $R_0>0$ and $\Lambda\ge1$. There are $\kappa\ge1$ and $c\ge1$ such that
for every $R>0$, every $x_0\in\R^d$ and every $f\colon B_{\kappa R}(x_0)\cap\Z^d\to\R$,
\[
  \sum_{\substack{x,y\in B_R(x_0)\cap\Z^d\\ \abs{x-y}>R_0}}(f(x)-f(y))^2\abs{x-y}^{-d-\alpha}
  \le c\sum_{\substack{x,y\in B_{\kappa R}(x_0)\cap\Z^d\\ \abs{x-y}>R_0}}(f(x)-f(y))^2\omega(x,y).
\]
The constant $c$ depends on $\Lambda$, $\vartheta$, $R_0$ and $d$. It does not depend on $\omega$
and $\Gamma$.
\end{bksthm}

The statement leaves the dependence of $\kappa$ open. The proof in \cite{BKS} gives
$\kappa=(N-1)\lambda$ (p.~29), and $\kappa$, like $c$, can be chosen depending only on $\Lambda$,
$\vartheta$, $R_0$ and $d$. We use Theorem~1.3 in this uniform form. So does the proof of Lemma~3.7
in \cite{BKS}, which fixes $\kappa$ and then applies Corollary~3.1 to configurations and kernels
that vary with $h$.

We use it with $R_0=\tfrac12$. Then the condition $\abs{x-y}>R_0$ holds for all distinct lattice
points, and the sums run over all pairs of distinct points, or all pairs, since diagonal terms
vanish.

\subsection{Random lattices}

Fix
\[
  \varepsilon=\min\bigl(\tfrac12,\tfrac12\sin(\vartheta/2)\bigr),\qquad \delta=\varepsilon/d .
\]
For $Q=(Q_1,\dots,Q_d)\in(\R^d)^d$ let $A_Q$ be the linear map $A_Qu=u+\sum_j u_jQ_j$, and let
$\Qd=\bar B_\delta(0)^d\subseteq(\R^d)^d$. For $Q\in\Qd$,
$\norm{A_Qu-u}\le\sum_j\abs{u_j}\delta\le d\delta\norm{u}=\varepsilon\norm{u}$. Hence
\begin{equation}\label{eq:nearid}
  (1-\varepsilon)\norm{u}\le\norm{A_Qu}\le(1+\varepsilon)\norm{u},\qquad
  (1-\varepsilon)^d\le\abs{\det A_Q}\le(1+\varepsilon)^d,\tag{$\ast$}
\end{equation}
the second because the singular values of $A_Q$ lie in $[1-\varepsilon,1+\varepsilon]$. Let
$P=[-\tfrac12,\tfrac12)^d$. For $h>0$, $Q\in\Qd$ and $\eta\in P$, the lattice map is
\[
  T(z)=T_{h,Q,\eta}(z)=h\,A_Q(z+\eta)\qquad(z\in\R^d),
\]
an injective affine map of $\R^d$, which we evaluate at the points of $\Z^d$.

\subsection{One lattice}

\begin{lemma}[Cones]\label{lem:cone}
Let $V$ be a double cone with axis $v$ and apex angle $\theta\ge\vartheta$, and let $V'$ be the
double cone with the same axis and apex angle $\vartheta/2$. Let $A$ be linear with
$\norm{Au-u}\le\varepsilon\norm{u}$ for all $u$, where $\varepsilon<\sin(\vartheta/2)$. Then
$AV'\subseteq V$.
\end{lemma}

\begin{proof}
Let $w\in V'$, so $w$ makes an angle less than $\vartheta/2$ with $v$ or with $-v$. Since
$\norm{Aw-w}\le\varepsilon\norm{w}$ and $\varepsilon<1$, $Aw\ne0$, and the angle between $Aw$ and
$w$ is at most $\arcsin\varepsilon<\vartheta/2$. By the triangle inequality for angles, $Aw$ makes
an angle less than $\vartheta\le\theta$ with the same half-line, so $Aw\in V$.
\end{proof}

\begin{lemma}[The inequality for one lattice]\label{lem:onelattice}
Let $\kappa_0$ and $c_0$ be the constants of the uniform form of Theorem~1.3 of \cite{BKS} for the parameters
$\vartheta/2$, $\Lambda'=2^{d+\alpha}\Lambda$ and $R_0=\tfrac12$. Let $\Gamma$ be
$\vartheta$-bounded, let $k$ be symmetric and satisfy \eqref{eq:2}, and let $f\colon\R^d\to\R$ be
any function. Write
\[
  F(s,t)=(f(t)-f(s))^2\norm{s-t}^{-d-\alpha},\qquad G(s,t)=(f(t)-f(s))^2k(s,t),
\]
and $F_D=F\1_{D\times D}$, $G_D=G\1_{D\times D}$. Then for every ball $B=B_R(x_0)$, every $h>0$,
$Q\in\Qd$ and $\eta\in P$, with $T=T_{h,Q,\eta}$ and $B^*=B_{3\kappa_0R}(x_0)$,
\[
  \sum_{z,z'\in\Z^d}F_B(Tz,Tz')\le 2^{d+\alpha}c_0\sum_{z,z'\in\Z^d}G_{B^*}(Tz,Tz').
\]
\end{lemma}

\begin{proof}
Write $A=A_Q$.

\emph{The pulled-back configuration.} Let $\Gamma'(z)$ be the double cone with the axis of
$\Gamma(Tz)$ and apex angle $\vartheta/2$, for every $z\in\R^d$. All its apex angles equal
$\vartheta/2$, so $\Gamma'$ is $\vartheta/2$-bounded. If
$z'\in V^{\Gamma'}[z]$, then $Tz'-Tz=hA(z'-z)\in\Gamma(Tz)$ by Lemma~\ref{lem:cone}, since double
cones are invariant under positive dilations. Hence
$\1_{V^{\Gamma'}[z]}(z')\le\1_{V^\Gamma[Tz]}(Tz')$.

\emph{The pulled-back kernel.} Let $\omega(z,z')=h^{d+\alpha}k(Tz,Tz')$, which is symmetric. By
\eqref{eq:nearid}, $h(1-\varepsilon)\norm{z-z'}\le\norm{Tz-Tz'}\le h(1+\varepsilon)\norm{z-z'}$.
With $\varepsilon\le\tfrac12$, \eqref{eq:2} gives, for $z\ne z'$,
\[
  \omega(z,z')\le\Lambda(1-\varepsilon)^{-d-\alpha}\norm{z-z'}^{-d-\alpha}
    \le\Lambda'\norm{z-z'}^{-d-\alpha},
\]
and
\[
  \omega(z,z')\ge\Lambda^{-1}(1+\varepsilon)^{-d-\alpha}
    \bigl(\1_{V^{\Gamma'}[z]}(z')+\1_{V^{\Gamma'}[z']}(z)\bigr)\norm{z-z'}^{-d-\alpha},
\]
with $\Lambda^{-1}(1+\varepsilon)^{-d-\alpha}\ge\Lambda'^{-1}$. These bounds hold at every pair
of lattice points, because \eqref{eq:2} holds at every pair of points; no measurability of
$\Gamma$ enters. So $(\Gamma',\omega)$ satisfies the hypotheses of Theorem~1.3 of \cite{BKS} with
$\vartheta/2$, $\Lambda'$ and $R_0=\tfrac12$.

\emph{The balls.} Let $\xi=T^{-1}x_0$ and $r=R/(h(1-\varepsilon))$. If $Tz\in B$ then
$\norm{z-\xi}\le\norm{Tz-x_0}/(h(1-\varepsilon))<r$, so $z\in B_r(\xi)$. If $z\in B_{\kappa_0r}(\xi)$
then $\norm{Tz-x_0}<h(1+\varepsilon)\kappa_0r=\kappa_0R(1+\varepsilon)/(1-\varepsilon)\le3\kappa_0R$,
so $Tz\in B^*$.

\emph{Conclusion.} Apply Theorem~1.3 of \cite{BKS} to $f\circ T$ on $B_r(\xi)$. Using
$\norm{Tz-Tz'}^{-d-\alpha}\le2^{d+\alpha}h^{-d-\alpha}\norm{z-z'}^{-d-\alpha}$ on the left,
\begin{align*}
  \sum_{z,z'}F_B(Tz,Tz')
  &\le2^{d+\alpha}h^{-d-\alpha}\sum_{z,z'\in B_r(\xi)}(f(Tz)-f(Tz'))^2\norm{z-z'}^{-d-\alpha}\\
  &\le2^{d+\alpha}c_0h^{-d-\alpha}\sum_{z,z'\in B_{\kappa_0r}(\xi)}(f(Tz)-f(Tz'))^2\omega(z,z').
\end{align*}
The last sum is $h^{d+\alpha}\sum_{z,z'\in B_{\kappa_0r}(\xi)}G(Tz,Tz')$, and every such pair has
$Tz,Tz'\in B^*$.
\end{proof}

\subsection{Averaging over lattices}

For measurable $H\colon\R^d\times\R^d\to[0,\infty]$ put
\[
  L_h(H)=\int_{\Qd}\int_P\sum_{z,z'\in\Z^d}H\bigl(T_{h,Q,\eta}z,\,T_{h,Q,\eta}z'\bigr)\dd\eta\dd Q .
\]
Here $\dd Q$ is Lebesgue measure on $(\R^d)^d$, and the integrand is jointly measurable in
$(Q,\eta)$.

\begin{lemma}[Tiling]\label{lem:tiling}
Let $H$ be measurable, nonnegative and zero on the diagonal. For $h>0$ and $Q\in\Qd$,
\[
  \int_P\sum_{z,z'}H(Tz,Tz')\dd\eta=\abs{\det(hA_Q)}^{-1}\sum_{w\in\Z^d\setminus\{0\}}
  \int_{\R^d}H(s,s+hA_Qw)\dd s .
\]
\end{lemma}

\begin{proof}
Write $z'=z+w$. The term $w=0$ vanishes. For fixed $w$, the translates $z+P$, $z\in\Z^d$, tile
$\R^d$, so $\sum_z\int_PH(hA(z+\eta),hA(z+\eta)+hAw)\dd\eta=\int_{\R^d}H(hAu,hAu+hAw)\dd u$.
The substitution $s=hAu$ finishes the proof. All terms are nonnegative, so sums and integrals
may be interchanged.
\end{proof}

\begin{lemma}[Averaging a column]\label{lem:column}
Let $g\colon\R^d\to[0,\infty]$ be measurable and $w\in\Z^d\setminus\{0\}$, with $m=\supn w$. Then
\[
  \abs{B_{\delta/d}}^{d-1}m^{-d}\int_{\bar B_{\delta m/d}(w)}g
  \;\le\;\int_{\Qd}g(A_Qw)\dd Q\;\le\;\abs{B_\delta}^{d-1}m^{-d}\int_{\bar B_{d\delta m}(w)}g .
\]
\end{lemma}

\begin{proof}
Choose $j$ with $\abs{w_j}=m$. Fix the columns $Q_i$, $i\ne j$, and put
$b=w+\sum_{i\ne j}w_iQ_i$, so that $A_Qw=b+w_jQ_j$. The substitution $y=b+w_jQ_j$ gives
$\int_{\bar B_\delta(0)}g(b+w_jQ_j)\dd Q_j=m^{-d}\int_{\bar B_{\delta m}(b)}g$.

For the upper bound, $\norm{b-w}\le(d-1)\delta m$, so $\bar B_{\delta m}(b)\subseteq\bar B_{d\delta m}(w)$.
Integrating over the other $d-1$ columns gives the factor $\abs{B_\delta}^{d-1}$.

For the lower bound, restrict the other columns to $\bar B_{\delta/d}(0)$. Then
$\norm{b-w}\le(d-1)\delta m/d$ and $\bar B_{\delta m}(b)\supseteq\bar B_{\delta m/d}(w)$, and the
restricted columns contribute $\abs{B_{\delta/d}}^{d-1}$.
\end{proof}

\begin{lemma}[Counting]\label{lem:count}
For $y\in\R^d$ put $\rho^+(y)=\sum_{w\ne0}\supn w^{-d}\1[\norm{y-w}\le d\delta\supn w]$ and
$\rho^-(y)=\sum_{w\ne0}\supn w^{-d}\1[\norm{y-w}\le\delta\supn w/d]$, the sums over
$w\in\Z^d\setminus\{0\}$. Then
\begin{enumerate}
\item $\rho^+(y)\le6^d$ for all $y$;
\item $\rho^-(y)\ge(\delta/(3d^2))^d$ whenever $\norm y\ge C_0:=2d^3/\delta$.
\end{enumerate}
\end{lemma}

\begin{proof}
(1) If $w$ contributes, then $\supn{y-w}\le\norm{y-w}\le d\delta\supn w\le\tfrac12\supn w$. Hence
$\tfrac12\supn w\le\supn y\le\tfrac32\supn w$. Since $\supn w\ge1$, this forces
$\supn y\ge\tfrac12$. The contributing $w$ lie in the cube $\{\supn{w-y}\le\supn y\}$, which
contains at most $(2\supn y+1)^d\le(4\supn y)^d$ lattice points. Each has weight
$\supn w^{-d}\le(\tfrac32/\supn y)^d$. So $\rho^+(y)\le6^d$.

(2) Let $M=\supn y\ge\norm y/\sqrt d\ge\norm y/d\ge2d^2/\delta$, and $r=\delta M/(2d^2)\ge1$.
Since $\delta\le1$, also $r\le M/2$. Consider $w\in\Z^d$ with $\supn{w-y}\le r$. Then
$\tfrac12M\le\supn w\le\tfrac32M$, in particular $w\ne0$, and
\[
  \norm{y-w}\le\sqrt d\,r\le d\,r=\delta M/(2d)\le\delta\supn w/d ,
\]
so $w$ contributes at least $(\tfrac32M)^{-d}$. Each coordinate interval
$[y_i-r,y_i+r]$ contains at least $\lfloor2r\rfloor\ge r$ integers, since $r\ge1$. So there are
at least $r^d$ such $w$, and $\rho^-(y)\ge r^d(\tfrac32M)^{-d}=(\delta/(3d^2))^d$.
\end{proof}

\begin{proposition}[The lattice average]\label{prop:avg}
Let $H\colon\R^d\times\R^d\to[0,\infty]$ be measurable and zero on the diagonal. For every $h>0$,
\[
  K_-h^{-2d}\iint_{\norm{t-s}\ge C_0h}H(s,t)\dd s\dd t\;\le\;L_h(H)\;\le\;K_+h^{-2d}\iint H(s,t)\dd s\dd t,
\]
with $K_+=6^d\abs{B_\delta}^{d-1}(1-\varepsilon)^{-d}$,
$K_-=(\delta/(3d^2))^d\abs{B_{\delta/d}}^{d-1}(1+\varepsilon)^{-d}$ and $C_0$ as in
Lemma~\ref{lem:count}.
\end{proposition}

\begin{proof}
By Lemma~\ref{lem:tiling}, \eqref{eq:nearid} and Tonelli's theorem,
\[
  L_h(H)\le h^{-d}(1-\varepsilon)^{-d}\sum_{w\ne0}\int_{\R^d}\int_{\Qd}H(s,s+hA_Qw)\dd Q\dd s .
\]
Apply Lemma~\ref{lem:column} with $g(y)=H(s,s+hy)$, and then sum over $w$:
\[
  L_h(H)\le h^{-d}(1-\varepsilon)^{-d}\abs{B_\delta}^{d-1}\iint H(s,s+hy)\,\rho^+(y)\dd y\dd s .
\]
By Lemma~\ref{lem:count}(1) and the substitution $t=s+hy$, for which
$\int H(s,s+hy)\dd y=h^{-d}\int H(s,t)\dd t$, this is at most $K_+h^{-2d}\iint H$.

The lower bound is the same computation, with the lower halves of \eqref{eq:nearid} and
Lemma~\ref{lem:column}. Discarding $\norm y<C_0$ and using Lemma~\ref{lem:count}(2) gives
\[
  L_h(H)\ge h^{-d}(1+\varepsilon)^{-d}\abs{B_{\delta/d}}^{d-1}\Bigl(\frac\delta{3d^2}\Bigr)^d
  \int_{\R^d}\int_{\norm y\ge C_0}H(s,s+hy)\dd y\dd s,
\]
and the substitution $t=s+hy$ turns $\{\norm y\ge C_0\}$ into $\{\norm{t-s}\ge C_0h\}$.
\end{proof}

\subsection{Conclusion}

\begin{proof}[Proof of Theorem~\ref{thm:A}]
With $\kappa_0,c_0$ as in Lemma~\ref{lem:onelattice},
put $\kappa=3\kappa_0$ and $B^*=B_{\kappa R}(x_0)$. Since $f$ and $k$ are measurable, $F_B$ and
$G_{B^*}$ are measurable, and both vanish on the diagonal, whatever the value of $k$ there. Lemma~\ref{lem:onelattice} holds for
every $(h,Q,\eta)$. Integrating it over $\eta\in P$ and $Q\in\Qd$ gives
$L_h(F_B)\le2^{d+\alpha}c_0L_h(G_{B^*})$. By Proposition~\ref{prop:avg},
\[
  K_-\iint_{\norm{t-s}\ge C_0h}F_B\le2^{d+\alpha}c_0K_+\iint G_{B^*}\qquad(h>0).
\]
As $h\downarrow0$ the sets $\{\norm{t-s}\ge C_0h\}$ increase to the complement of the diagonal,
a null set. By monotone convergence,
\[
  \abs{f}^2_{H^{\alpha/2}(B)}=\iint F_B\le\frac{2^{d+\alpha}c_0K_+}{K_-}\iint G_{B^*}
  =C\,\abs{f}^2_{H_k(B^*)} .
\]
The constants $\varepsilon,\delta,K_\pm,C_0$ depend only on $d$ and $\vartheta$, and $\kappa_0,c_0$
only on $d$, $\vartheta$, $\Lambda$ and $\alpha$.
\end{proof}

\begin{remark}\label{rem:alpha}
The constants $\kappa$ and $C$ of Theorem~\ref{thm:A} can be chosen independently of
$\alpha\in(0,2)$. Indeed, so can the constants of the uniform form of Theorem~1.3
(Section~\ref{sec:discrete}). Its proof produces a factor $\lambda^{d+\alpha}$, where $\lambda\ge R_0$
is the constant of Theorem~5.15 of \cite{BKS}. Since $\lambda$ may be taken $\ge1$, this is at most
$\lambda^{d+2}$. Then use $\Lambda'=2^{d+2}\Lambda$ in Lemma~\ref{lem:onelattice}, and
$2^{d+\alpha}\le2^{d+2}$. Both uniform statements are formalized (Section~\ref{sec:lean}).
\end{remark}

\section{Consequences}\label{sec:consequences}

A configuration is \emph{$\vartheta$-admissible} in \cite{BKS} if it is $\vartheta$-bounded and
$\{(x,y):y-x\in\Gamma(x)\}$ is Borel (condition (M)). Theorem~\ref{thm:A} does not use (M).

\begin{corollary}\label{cor:bks}
Lemma~3.7, Theorem~1.1, and the comparability statements of Theorem~1.4 of \cite{BKS} hold as
stated.
\end{corollary}

\begin{proof}
In the proof of Lemma~3.7, the passage in Section~\ref{sec:gap} is meant to establish
$c\,\abs{f}^2_{H^{\alpha/2}(B)}\le\abs{f}^2_{H_k(B^*)}$ for all $f\in H_k(B^*)$ and every ball, with
$c$ independent of the ball. Theorem~\ref{thm:A} gives this for every $f\in L^2(B^*)$: take a
measurable representative, extended by zero, and note that both sides depend only on the values of
$f$ on $B^*$. The rest of the proof of Lemma~3.7 in \cite{BKS} is the application of their
Lemma~A.1, which removes the enlargement and is unaffected by the gap. Theorem~1.1 follows from
Lemma~3.7 as in \cite{BKS}. So do the comparabilities in Theorem~1.4: on $\R^d$ by letting the
radius tend to infinity, and on bounded Lipschitz domains through Lemma~A.1.
\end{proof}

The dependence on $\alpha_0$ asserted in Theorem~1.1 follows from Remark~\ref{rem:alpha} and the
corresponding statement of Lemma~A.1.

\subsection{The route through Chaker and Silvestre}

Chaker and Silvestre \cite[Assumption~1.1]{CS} ask for $\mu\in(0,1)$ and $\lambda>0$ such that
for every ball $B$ and every $x\in B$,
$\abs{\{z\in B:K(x,z)\ge\lambda\norm{x-z}^{-d-2s}\}}\ge\mu\abs B$.

\begin{lemma}\label{lem:cs}
Let $\Gamma$ be $\vartheta$-bounded, and let $k$ be symmetric and satisfy \eqref{eq:2}. Then for
every ball $B=B_r(p)$ and every $x\in B$,
\[
  \abs{\{z\in B:k(x,z)\ge\Lambda^{-1}\norm{x-z}^{-d-\alpha}\}}\ge\Bigl(\frac{\sin^2\vartheta}{16}\Bigr)^d\abs B .
\]
That is, $k$ satisfies Assumption~1.1 of \cite{CS} with $s=\alpha/2$, $\lambda=\Lambda^{-1}$ and
$\mu=(\sin^2\vartheta/16)^d$.
\end{lemma}

\begin{proof}
Put $\sigma=\sin\vartheta$. Let $v$ be the axis of $\Gamma(x)$, with its sign chosen so that
$\langle v,p-x\rangle\ge0$. Since the apex angle of $\Gamma(x)$ is at least $\vartheta$, $\Gamma(x)$
contains the half-cone $W=\{h\ne0:\angle(h,v)<\vartheta\}$. For $h\in\R^d$ put
$g(h)=\langle v,h\rangle\sigma-\norm{h-\langle v,h\rangle v}\cos\vartheta$. Then
$g(h)=\norm h\sin(\vartheta-\angle(h,v))$ for $h\ne0$, so $g>0$ exactly on $W$. By the
Cauchy--Schwarz inequality for the pair $(\sigma,\cos\vartheta)$, $g$ is $1$-Lipschitz, and
$g(h+tv)=g(h)+t\sigma$.

Let $a=r\sigma/4$, $\varepsilon=\sigma^2/8$, $\rho=r\sigma^2/16$ and
$q=x+av+\varepsilon(p-x)$. We show that $B_\rho(q)\subseteq B$ and $B_\rho(q)-x\subseteq W$.

First, $q-p=av-(1-\varepsilon)(p-x)$, and $\langle v,p-x\rangle\ge0$, $\norm{p-x}<r$, so
\[
  \norm{q-p}^2\le a^2+(1-\varepsilon)^2r^2=(r-\rho)^2-\frac{r^2\sigma^2(16-3\sigma^2)}{256}\le(r-\rho)^2 .
\]
Hence $\norm{q-p}+\rho\le r$, and $B_\rho(q)\subseteq B$.

Second, $g(q-x)=g(\varepsilon(p-x))+a\sigma\ge-\varepsilon\norm{p-x}+a\sigma>-\varepsilon r+a\sigma
=r\sigma^2/8$, using $g(0)=0$. For $z\in B_\rho(q)$, $g(z-x)\ge g(q-x)-\norm{z-q}>r\sigma^2/8-\rho>0$,
so $z-x\in W\subseteq\Gamma(x)$.

For $z\in B_\rho(q)$, therefore, $z\in B$ and $z\in V^\Gamma[x]$, and \eqref{eq:2} gives
$k(x,z)\ge\Lambda^{-1}\norm{x-z}^{-d-\alpha}$. Finally $\abs{B_\rho(q)}=(\rho/r)^d\abs B
=(\sigma^2/16)^d\abs B$.
\end{proof}

\begin{remark}\label{rem:cs}
By Lemma~\ref{lem:cs}, \cite[Theorem~1.3]{CS} applies to every kernel satisfying \eqref{eq:2}.
Rescaling $B_R(x_0)$ to $B_1(0)$ preserves \eqref{eq:2}, since double cones are invariant under
positive dilations, so it gives the enlarged-ball inequality of Theorem~\ref{thm:A} with
$\kappa=2$.
\end{remark}

\section{Formalization}\label{sec:lean}

The argument above is checked in Lean~4 with Mathlib, as part of a formalization of \cite{BKS}.
The repository is \url{https://github.com/dbenbenn/quadratic-forms-sobolev-lean}. The headline
results are stated in Lean exactly as above, with $\vartheta$-boundedness as in Definition~2.1 of
\cite{BKS} (\texttt{QFS.IsThetaBounded}):

\begin{center}
\begin{tabular}{ll}
Theorem~\ref{thm:A} & \texttt{QFS.GapNote.theorem\_A}\\
Remark~\ref{rem:alpha} & \texttt{QFS.GapNote.remark\_7}\\
Lemma~\ref{lem:cs} & \texttt{QFS.GapNote.lemma\_9}\\
Theorem~1.3 of \cite{BKS} & \texttt{QFS.Paper.theorem\_1\_3}
\end{tabular}
\end{center}

The first three are in the file \texttt{GapNote.lean}, the last in \texttt{Paper.lean}, and the
lemmas in the directory \texttt{RandomLattice}, all under \texttt{QuadraticFormsSobolev}:

\begin{center}
\begin{tabular}{ll}
Lemma~\ref{lem:cone} & \texttt{QFS.mem\_doubleCone\_of\_near}\\
Lemma~\ref{lem:onelattice} & \texttt{QFS.latticeSum\_ball\_le}\\
Lemma~\ref{lem:tiling} & \texttt{QFS.lintegral\_latticeSum\_le}, \texttt{QFS.le\_lintegral\_latticeSum}\\
Lemma~\ref{lem:column} & \texttt{QFS.lintegral\_colBox\_le}, \texttt{QFS.le\_lintegral\_colBox}\\
Lemma~\ref{lem:count} & \texttt{QFS.tsum\_count\_le}, \texttt{QFS.le\_tsum\_count}\\
Proposition~\ref{prop:avg} & \texttt{QFS.latAvg\_le}, \texttt{QFS.le\_latAvg}\\
Uniform Theorem~1.3 & \texttt{QFS.theoremOneThree\_uniform}
\end{tabular}
\end{center}

These are the forms the Lean proof uses, not transcriptions. The Lean form of
Lemma~\ref{lem:onelattice} assumes only that $\vartheta$ is a lower bound for the apex angles
(\texttt{QFS.ApexLowerBound}), and allows $\alpha\ge0$. The two tiling inequalities combine Lemma~\ref{lem:tiling} with the bound
\eqref{eq:nearid} on the determinant. Some constants differ harmlessly: $4^d$ in
Lemma~\ref{lem:count}(1) and $(\delta/(4d^2))^d$ in Lemma~\ref{lem:count}(2).

The example of Section~\ref{sec:gap}, where the limit step of \cite{BKS} fails for $\alpha>1$, is
formalized on $\R$ in the file \texttt{BoundaryDefect.lean}.

The formalization also proves Theorem~1.3 of \cite{BKS} in full, in the $\alpha$-uniform form of
Remark~\ref{rem:alpha}, and Theorem~1.1 with the $\alpha_0$ clause. The statement of
Theorem~1.3 is correct as printed. A few steps of its printed proof needed local repairs in the
formalization; they are independent of the gap discussed here and are not treated in this note. For the step removing the
enlargement, it gives a proof of Lemma~A.1 that avoids the Whitney decomposition, with Dyda's
inequality proved rather than quoted. All results depend only on the standard axioms
(\texttt{propext}, \texttt{Classical.choice}, \texttt{Quot.sound}).

\begin{thebibliography}{9}
\bibitem{BKS} K.-U. Bux, M. Kassmann, and T. Schulze, \emph{Quadratic forms and Sobolev spaces of
fractional order}, Proc. London Math. Soc.\ (3) \textbf{119} (2019), 841--866,
\href{https://doi.org/10.1112/plms.12246}{doi:10.1112/plms.12246}. Cited from
\href{https://arxiv.org/abs/1707.09277v1}{arXiv:1707.09277v1}, whose numbering and page numbers we
use; we were unable to consult the journal version.
\bibitem{CS} J. Chaker and L. Silvestre, \emph{Coercivity estimates for integro-differential
operators}, Calc. Var. Partial Differential Equations \textbf{59} (2020), Paper No.~106,
\href{https://doi.org/10.1007/s00526-020-01764-y}{doi:10.1007/s00526-020-01764-y},
\href{https://arxiv.org/abs/1904.13014}{arXiv:1904.13014}.
\bibitem{Dyda} B. Dyda, \emph{On comparability of integral forms}, J. Math. Anal. Appl.\
\textbf{318} (2006), 564--577,
\href{https://doi.org/10.1016/j.jmaa.2005.06.021}{doi:10.1016/j.jmaa.2005.06.021}.
\end{thebibliography}

\bigskip
\noindent\textit{About this note.} This note was written by Claude, an AI model made by
Anthropic. The mathematics in it was developed and machine-checked in the Lean formalization
cited in Section~\ref{sec:lean}.

\end{document}
